% Generated draft: requires mathematical and source review.
\documentclass[11pt]{article}
\usepackage[margin=25mm]{geometry}
\usepackage{fontspec}
\setmainfont{Latin Modern Roman}
\usepackage{amsmath,amssymb,mathtools,unicode-math}
\setmathfont{Latin Modern Math}
\usepackage{enumitem,needspace,xurl,hyperref}
\setlength{\parindent}{0pt}
\setlength{\parskip}{5pt}
\setlength{\emergencystretch}{3em}
\newcommand{\N}{\mathbb{N}}
\newcommand{\Z}{\mathbb{Z}}
\newcommand{\Q}{\mathbb{Q}}
\newcommand{\R}{\mathbb{R}}
\newcommand{\C}{\mathbb{C}}
\newcommand{\F}{\mathbb{F}}
\newcommand{\A}{\mathbb{A}}
\newcommand{\PP}{\mathbb{P}}
\newcommand{\E}{\mathbb{E}}
\newcommand{\eps}{\varepsilon}
\newcommand{\dd}{\,\mathrm{d}}
\begin{document}
\section*{UnsolvedMath ID: 6800008}
\textbf{Source problem code:} catalogue entry
\subsection{Definitions and mathematical statement}
Let $n\geq2$ and $1\leq k\leq2n$, so $2n+1>3$ and $0<k<2n+1$. For a real $(k+1)$-plane $V\subset\mathbb R^{2n+2}$, the equatorial sphere $S(V)=V\cap S^{2n+1}$ projects under the antipodal quotient to a projective subspace $Q_V\cong\mathbb RP^k$ in $\mathbb RP^{2n+1}$.

The source equips the ambient projective space with a Berger metric $g_\mu$ without giving its normalization or allowed parameter range. A usual Riemannian convention on the covering sphere is $g_\mu=g_{\mathrm{round}}+(\mu-1)\eta\otimes\eta$, $\mu>0$, where $\eta$ is dual to the unit Hopf vector field $z\mapsto iz$; this descends under the antipodal quotient. This formula is explanatory and must be checked against the intended source convention.

Are the projected equatorial subspaces under consideration minimal for $g_\mu$, meaning that their mean-curvature vector vanishes? Are they volume minimizing in their respective homology classes, and are all volume-minimizing cycles in those classes of this projective-subspace form?

For example, with $\mathbb Z/2$ coefficients, the minimization comparison would be
\[
\mathbf M_{g_\mu}(T)\geq\operatorname{Vol}_{g_\mu}(Q_V)
\quad\text{whenever }[T]=[Q_V]\in H_k(\mathbb RP^{2n+1};\mathbb Z/2),
\]
followed by classification of all equality cases. The source does not fix coefficients or cycle regularity. This matters because even-dimensional real projective spaces are nonorientable. It also does not specify whether every real plane $V$ is included or a distinguished family relative to the Hopf structure. The claim that only projective-subspace cycles minimize is not the assertion that a single representative is unique within its homology class.
\subsection{Short English statement}
For Berger metrics on odd-dimensional real projective space, determine whether projected equatorial spheres are minimal and homologically volume minimizing, and whether every minimizer has that form. Fix the metric, subspace family, and homology coefficients.
\subsection{Sources}
\begin{itemize}
\item Morgan and Pansu, \emph{A List of Open Problems in Differential Geometry} (2018), Question environment 8, ``Area minimizing projective spaces in the projective space with the Berger metric,'' proposed by Olga Gil-Medrano, University of Valencia, Spain, \url{https://www.imo.universite-paris-saclay.fr/~pansu/problems_MTDG.tex}; accessed 2026-07-29, source-TeX extraction.
\item Local registry \texttt{data/unsolvedmath/problems.json}, permanent UnsolvedMath ID $6800008$, source code catalogue entry, same titled problem. Dataset attribution: Ulam AI and the UnsolvedMath Contributors.
\end{itemize}
\end{document}
